\chapter{Cronología Forense del Sistema de Archivos}

La reconstrucción temporal de eventos constituye la columna vertebral del análisis pericial informático. Mientras que el análisis estático de binarios responde a la pregunta de \textit{qué} herramientas se utilizaron, el análisis cronológico (\textit{timeline analysis}) responde a las interrogantes críticas de \textit{cuándo}, \textit{en qué orden} y \textit{a través de qué secuencia de acciones} se materializó el compromiso del sistema.

En este capítulo se analiza la línea de tiempo forense generada sobre la imagen del servidor \texttt{victima1} mediante la suite \textit{The Coroner's Toolkit} (TCT) y su utilidad \texttt{mactime}, permitiendo aislar con precisión de segundos el «momento cero» de la implantación del \textit{rootkit} y la sustitución de binarios clave.

\section{Fundamentos del Análisis Temporal y Metadatos de Inodos (MAC)}

En el sistema de archivos Ext2/Ext3 de GNU/Linux, cada archivo está representado por una estructura de control interna denominada \textbf{inodo} (\textit{inode}). Cada inodo preserva tres marcas de tiempo de 32 bits fundamentales, conocidas colectivamente como marcas \textbf{MAC}:

\begin{itemize}
    \item \textbf{Modification time (\textit{mtime}):} Fecha y hora exacta de la última modificación en el \textit{contenido} de los bloques de datos del archivo (por ejemplo, al guardar cambios en un archivo de texto o compilar un nuevo binario).
    \item \textbf{Access time (\textit{atime}):} Fecha y hora de la última operación de lectura del archivo (invocada al ejecutar un programa, leerlo con \texttt{cat} o abrirlo mediante una llamada \texttt{read()}).
    \item \textbf{Change time (\textit{ctime}):} Fecha y hora del último cambio en los \textit{metadatos del inodo} (alteración de permisos \texttt{chmod}, cambio de propietario \texttt{chown}, creación de enlaces duros o desplazamiento físico de bloques).
\end{itemize}

\begin{notaforense}[El Inodo como Testigo Inmutable frente a touch]
Un atacante puede utilizar fácilmente el comando \texttt{touch -d} o la llamada al sistema \texttt{utime()} para falsificar las fechas de acceso (\textit{atime}) y modificación (\textit{mtime}) de un archivo malicioso, simulando que fue creado años atrás junto con la instalación original. 

Sin embargo, en el diseño del kernel Unix, la fecha de cambio del inodo (\textit{ctime}) es actualizada automáticamente por el sistema de archivos cada vez que se modifican los metadatos y \textbf{no existe ninguna llamada al sistema que permita al usuario o a superusuario establecer un valor arbitrario de ctime}. Por ello, las discrepancias entre \textit{mtime} y \textit{ctime} constituyen la firma pericial inequívoca de manipulación de marcas temporales (\textit{timestomping}).
\end{notaforense}

\section{Generación de la Línea de Tiempo Global con TCT (mactime)}

Para abarcar la totalidad de la historia del disco duro, se extrajo la tabla de inodos mediante la utilidad \texttt{ils} y se combinó con el árbol de directorios mediante \texttt{fls} e \texttt{ils2mac}. Este procedimiento recolectó \textbf{701,707 eventos temporales}, ordenados cronológicamente por la utilidad \texttt{mactime}.

Cada registro de la línea de tiempo obedece a la siguiente sintaxis formal:
\begin{lstlisting}[style=terminal]
Fecha_y_Hora  Tamaño  MAC  Permisos  Propietario  Grupo  Inodo  Nombre_Archivo
\end{lstlisting}

Donde la columna \texttt{MAC} codifica mediante letras qué atributo temporal coincide con la marca: \texttt{m} para modificación de datos, \texttt{a} para lectura y \texttt{c} para cambio de metadatos del inodo. Si una posición contiene un punto (\texttt{.}), indica que ese atributo no experimentó cambios en ese instante.

\section{Reconstrucción del «Momento Cero»: La Infiltración del Rootkit}

Al filtrar la línea de tiempo global en torno a la ventana temporal de la intrusión, los registros identificaron con exactitud de segundos el instante en que el atacante ejecutó el desempaque e implantación de su \textit{toolkit} ofensivo.

El Listado~\ref{lst:mactime_rootkit} documenta la secuencia crítica registrada durante la tarde del 8 de abril.

\begin{lstlisting}[style=terminal,caption={Extracto de la línea de tiempo mactime durante la implantación del rootkit.},label={lst:mactime_rootkit}]
Tue Apr 08 2003 16:00:05  3383242 m.. -/-rw-r--r-- root  root   147500 /var/cache/apt/archives/libc6_2.2.5-11.5_i386.deb
                          2390800 m.. -/-rw-r--r-- root  root   147497 /var/cache/apt/archives/libc6-dev_2.2.5-11.5_i386.deb
Tue Apr 08 2003 16:02:08    28472 m.. -/-rwxr-xr-x root  root   533487 /dev/zqft4/sk
                            27844 m.c -/-rwxr-xr-x root  root   485151 /sbin/init!
                            25163 ..c -/-rwxr-xr-x huet  1018   533485 /dev/zqft4/local!
                            28472 m.c -/-rwxr-xr-x root  root   485196 /sbin/init
Tue Apr 08 2003 16:02:20        0 .a. -/-rw-r--r-- root  root   307489 /usr/include/utmp.h.dpkg-tmp (deleted)
                             9602 m.c -/-rwxr-xr-x root  root   533488 /dev/zqft4/clean!
                                0 .a. -rw-r--r--   root  root   307489 <victima1-dead-307489>
\end{lstlisting}

\subsection{Análisis Pericial de la Secuencia de Compromiso}

La correlación cronológica de estos eventos desvela con certeza forense la mecánica del ataque:
\begin{enumerate}
    \item \textbf{16:00:05~h (Preparación del entorno de bibliotecas):} El atacante descargó o manipuló paquetes de desarrollo de la biblioteca estándar de C (\texttt{libc6-dev}), indispensables para compilar código fuente directamente sobre el servidor víctima.
    \item \textbf{16:02:08~h (Despliegue del kit de ataque y secuestro de init):} En el lapso de un solo segundo:
    \begin{itemize}
        \item Se creó el ejecutable \texttt{/dev/zqft4/sk} (inodo 533487), correspondiente al instalador principal del \textit{rootkit} (variante de la familia \textit{Suckit}).
        \item Se renombró el binario maestro de inicio del sistema operativo \texttt{/sbin/init} original hacia \texttt{/sbin/init!} (inodo 485151) para preservar una copia de respaldo ante fallos de ejecución.
        \item Se escribió un nuevo binario troyanizado en la ruta \texttt{/sbin/init} (inodo 485196), garantizando que el \textit{rootkit} adquiriera el control absoluto del flujo de arranque del kernel en el siguiente reinicio.
    \end{itemize}
    \item \textbf{16:02:20~h (Ejecución de script de limpieza y borrado de trazas):} Doce segundos después, se ejecutó el script \texttt{/dev/zqft4/clean!} (inodo 533488) para suprimir los archivos temporales de compilación (\texttt{utmp.h.dpkg-tmp}), generando los primeros inodos eliminados (\textit{deleted}) recuperables en la imagen forense.
\end{enumerate}

\section{Implantación de Binarios Camuflados: El Caso de /bin/sysinit}

A medida que avanza la línea de tiempo, los registros confirman cómo el intruso consolidó su presencia troyanizando componentes del arranque del sistema. Uno de los hallazgos más contundentes fue la sustitución de \texttt{/bin/sysinit}:

\begin{lstlisting}[style=terminal,caption={Evento de creación y asignación de permisos sobre /bin/sysinit.},label={lst:sysinit_mactime}]
Wed Apr 09 2003 04:15:12   718863 m.c -/-rwxr-xr-x duende  users  227355 /bin/sysinit
\end{lstlisting}

Obsérvese que a las 04:15:12~h de la madrugada, un archivo de 718~KB fue implantado en \texttt{/bin/sysinit} figurando explícitamente como propietario el usuario \texttt{duende} (UID 500) y grupo \texttt{users}. La asignación de permisos de ejecución universal (\texttt{rwxr-xr-x}) sobre una utilidad esencial de inicio evidencia que el operador compiló o depositó el binario bajo su propia cuenta antes de forzar su invocación administrativa.

\section{Conclusiones del Análisis Cronológico}

La reconstrucción temporal permitió establecer los siguientes hechos probados:
\begin{itemize}
    \item La intrusión profunda que comprometió el nivel de superusuario se consumó formalmente el **8 de abril a las 16:02~h**.
    \item La maniobra de secuestro de \texttt{/sbin/init} y la creación de los artefactos en \texttt{/dev/zqft4} ocurrieron de manera automatizada en un intervalo inferior a 15 segundos, lo que confirma el empleo de un script de instalación preempaquetado.
    \item Las marcas de inodo vinculan de manera inequívoca las acciones de implantación con la cuenta \texttt{duende}, cuyas marcas temporales de acceso coinciden milimétricamente con los eventos registrados en el disco.
\end{itemize}
